\section*{15. 受限最短路径（带钱数约束）}

给定无向图$G=(V,E)$，边$e_{uv}$长度为$l_{uv}$，顶点$u$花费为$c[u]$。从$s$出发到$t$，初始钱数为$M$。每经过顶点$u$需支付$c[u]$（剩余钱数$\ge c[u]$才能通过），求$s$到$t$的最短路径长度。



\begin{itemize}
    \item \textbf{状态扩展}：$dist[u][m]$表示到达顶点$u$、剩余钱数为$m$时的最短路径长度
    \item \textbf{状态转移}：对边$(u,v)$，$$dist[v][m-c[v]] = \min(dist[v][m-c[v]], dist[u][m] + l_{uv})$$
    \item \textbf{初始化}：$dist[s][M] = 0$，其余为$\infty$
    \item \textbf{优先队列}：按路径长度从小到大处理状态
\end{itemize}
\begin{lstlisting}[language=C++, style=cppstyle]
struct Edge {
    int to, len;
};

struct State {
    int u, money;
    int dist;
    bool operator>(const State& other) const {
        return dist > other.dist;
    }
};

int constrainedShortestPath(int n, int M, int s, int t,
                            vector<int>& cost,
                            vector<vector<Edge>>& adj) {
    // dist[u][m] = 到达u、剩余m元的最短距离
    vector<vector<int>> dist(n, vector<int>(M + 1, INT_MAX / 2));
    dist[s][M] = 0;

    // 优先队列（小根堆）
    priority_queue<State, vector<State>, greater<State>> pq;
    pq.push({s, M, 0});

    while (!pq.empty()) {
        State cur = pq.top(); pq.pop();

        int u = cur.u, m = cur.money, d = cur.dist;

        if (d > dist[u][m]) continue;  // 已优化
        if (u == t) return d;  // 到达终点

        for (Edge& e : adj[u]) {
            int v = e.to;
            int newMoney = m - cost[v];

            if (newMoney >= 0 && d + e.len < dist[v][newMoney]) {
                dist[v][newMoney] = d + e.len;
                pq.push({v, newMoney, dist[v][newMoney]});
            }
        }
    }

    return -1;  // 不可达
}
\end{lstlisting}


\begin{enumerate}
    \item 状态扩展：dist[u][m]表示"到达u且剩余钱为m时的最短距离"
    \item 起点dist[s][M]=0，其余为无穷大
    \item 用优先队列按距离处理，对出边松弛操作
\end{enumerate}

